\section{Conclusion}
\label{sec:conclusion}

In summary, we have presented our initial efforts towards analyzing Cornerstone ~\cite{cornerstone}. We developed a harness for running the application using GPUs or CPUs and with either a local octree setup or a distributed setup. Around this we created scripts to run parameter sweeps and gather profiles and the octree view. We also developed scripts to generate the distributions of initial positions and their perturbations. The Cornerstone ~\cite{cornerstone} code and ours were instrumented using the NVTX profiling tools and PCL (perf). 

To assess the update time after applying a perturbation, we used a CloudLab node with V100 GPUs. We showed that perturbation strength did correlate with the rebalancing time but only with larger numbers of particles. On the A100s, we determined the sorting and reordering steps dominate the computation of the tree locally, and were able to make significant optimizations to such construction. We also showed that the distributed case introduced significant overhead near two orders of magnitude but was not due to communication costs, which needs further exploration.
